% Source: 64_Notes.txt
\backchapter[Notes]{NOTES}{bm:notes}

\noteshead{Chapter 1}
\begin{notelist}
\entry The opening description of crucifixion reflects Bamber Gascoigne, The Christians (New York:
  Morrow, 1977), 17.
\entry The quotations in this chapter, as well as many vivid details, are from the descriptions of
  Jesus’ times in Great People of the Bible and How They Lived (Pleasantville: Reader’s Digest
  Association. 1968), 308, 338, 370, 379–81.
\end{notelist}

\noteshead{Chapter 2}
\begin{notelist}
\entry The description of Antioch as well as the quotation about Jerusalem and its fall are from
  Great People of the Bible, 406, 407, and 390.
\end{notelist}

\noteshead{Chapter 3}
\begin{notelist}
\entry For the text box see Rodney Stark, The Rise of Christianity: A Sociologist Reconsiders
  History. (Princeton: Princeton University Press, 1996).
\entry The story about King Abgar is found in Eusebius’s Ecclesiastical History, book I, section
  XIII.
\entry Professor Ward Gasque’s introduction to Irenaeus and Tertullian is found in Eerdman’s
  Handbook to the History of Christianity (Grand Rapids: Eerdmans, 1977), 75–77. This volume, edited
  by Tim Dowley, is one of the best introductions to church history available.
\entry Celsus’ criticism of Christians comes from Origen’s Against Celsus, book III, section 44.
\entry Julian’s views of Christian behavior are in a letter to Arsacius, high priest of Galatia,
  recorded in Sozomen’s Ecclesiastical History V. 16.
\end{notelist}

\noteshead{Chapter 4}
\begin{notelist}
\entry The Martyrdom of Polycarp is one of the better known stories from the early church. It can be
  found in Early Christian Fathers, edited by Cyril C. Richardson (Philadelphia: Westminster, 1953),
  141–58.
\entry The contrast between Jewish and Christian proselyting I have drawn from Paul Hutchinson and
  Winfred E. Garrison, 20 Centuries of Christianity (New York: Harcourt, Brace and Co., 1959),
  30–31.
\entry Pliny’s letter to Trajan is recorded in A New Eusebius, edited by J. Stevenson (London:
  S.P.C.K., 1960), 13–14.
\entry I found The Early Church by Henry Chadwick (Middlesex: Penguin, 1967), 54–60, helpful in
  “Reasons for the Gospel’s Spread.”
\end{notelist}

\noteshead{Chapter 5}
\begin{notelist}
\entry The Mahatma Gandhi quotation is echoed in Robert Johnson, The Meaning of Christ
  (Philadelphia: Westminster, 1958), 63.
\entry The discussion on “Faith and Theology” reflects the insights of J. W. C. Wand, The Four Great
  Heresies (London: A. R. Mowbray, 1955), chapter 1.
\entry The Charles Bigg quote can be found in Charles Williams, Descent of the Dove (New York:
  Meridian Books, 1956), 23.
\entry The Old Roman Creed appears in Documents of the Christian Church, edited by Henry Bettenson
  (London: Oxford University Press, 1963), 23–24.
\entry The C. S. Lewis quotation is from his Mere Christianity (New York: Macmillan, 1952), book 2,
  chapter 5.
\entry The description of orthodoxy under “Faith and History” and the comparison of virgin and born
  are from William Hordern, A Layman’s Guide to Protestant Theology (New York: Macmillan, 1974), 1,
  13.
\entry The text box “What Gnostics Believe” is adapted from J.N.D. Kelley, Early Christian
  Doctrines, 26.
\entry The text box “The Gospel of Thomas,” is adapted from Craig Evans, Fabricating Jesus (Downers
  Grove, IL: InterVarsity Press, 2006), 71.
\entry The imagery found in text box “Understanding Gnosticism Today” arose out of a lecture titled
  “Did We Get Jesus Right? Jesus in the Canonized and Apocryphal Gospels,” by Simon Gathercole and
  response by David Chapman at Lanier Theological Library on September 8, 2012.
\end{notelist}

\noteshead{Chapter 6}
\begin{notelist}
\entry The account of the believer in Sicily comes from Herbert B. Workman, Persecution in the Early
  Church (London: Charles H. Kelly, 1906), 275–76.
\entry Origen’s views about the meanings of Scripture come from his work On First Principles, Book
  IV, sections 7–8.
\entry “The Question of the Apocrypha” reflects the clear discussion of Floyd Filson, Which Books
  Belong in the Bible? (Philadelphia: Westminster, 1957).
\end{notelist}

\noteshead{Chapter 7}
\begin{notelist}
\entry The quotations from “One Early Christian” and from Athenagoras can be found in Adolph
  Harnack, Mission and Expansion of Early Christianity (New York: Harper, 1962), 207–9.
\entry This chapter draws upon the discussion “Forgiveness of Sins” in Roland H. Bainton,
  Christendom, vol. 1 (New York: Harper \& Row, 1964, 1966). I have reshaped Bainton’s helpful
  explanation of Cyprian’s attitude.
\entry Adapted from Everett Ferguson, Church History, Volume One: From Christ to the Pre-Reformation
  (Grand Rapids: Zondervan Press, 2005), 114–119.
\end{notelist}

\noteshead{Chapter 8}
\begin{notelist}
\entry Tertullian’s tirade against philosophy is from his work Prescriptions Against the Heretics,
  section 7. The quotation by Origen about the purpose of creation is from his work On First
  Principles, book II, section II, 4.
\entry The quotation from Clement about wisdom comes from A History of the Church by August Franzen
  (New York: Herder and Herder, 1969), 37.
\end{notelist}

\noteshead{Chapter 9}
\begin{notelist}
\entry The description of Diocletian and Galerius is a reshaped account from Hutchinson and
  Garrison, 20 Centuries of Christianity, 44–48. The quotations are from this work.
\entry Ambrose’s letter to Theodosius, letter 51, is quoted in Robert Payne, Fathers of the Western
  Church (New York: Viking Press, 1951), 78–79.
\entry The quotation by Bamber Gascoigne is from his The Christians, 44–45.
\end{notelist}

\noteshead{Chapter 10}
\begin{notelist}
\entry The quotation from the bishop is from Gregory of Nyssa and is quoted in W. H. C. Frend, The
  Early Church (Philadelphia: Lippincott, 1966), 186–87.
\entry The Nicene Creed appears, among many other places, in Philip Schaff, Creeds of Christendom,
  vol. 2 (New York: Harper, 1919), 58–59.
\entry Eusebius of Caesarea’s enthusiastic remarks are from his Life of Constantine, III, 15.
\entry William Hordern’s anecdote is in his book A Layman’s Guide to Protestant Theology,15–16.
\entry In the discussion of the Trinity, I found help in Fisher Humphrey, The Almighty (Elgin, IL:
  David C. Cook, 1976), esp. 102–7; Compare Roger Olson, The Story of Christian Theology: Twenty
  Centuries of Tradition \& Reform (Downers Grove, IL: InterVarsity Press, 1999); also see also
  Justo Gonzalez, A Concise History of Christian Doctrine (Nashville, TN: Abingdon, 2006).
\end{notelist}

\noteshead{Chapter 11}
\begin{notelist}
\entry J. S. Whale’s observation at Cambridge is from his book Christian Doctrine (London: Fontana
  Books, 1957), 102.
\entry David F. Wright’s discussion of the Christology debate is in Eerdmans’ Handbook to the
  History of Christianity, 156ff. The quotation is on page 171.
\entry The Chalcedonian definition may be found in Henry Bettenson, Documents of the Christian
  Church, 51–52.
\entry The chart is drawn from numerous handbooks on doctrine. One can see elements and influence
  adapted from Strong, Gonzalez, and Coppedge. Allen Coppedge, The God Who is Triune: Revisioning
  the Christian Doctrine of God, (Downers Grove, IL: InterVarsity, 2007); Justo Gonzalez, A Concise
  History of Christian Doctrine (Nashville, TN: Abingdon, 2006); and A. H. Strong’s, Systematic
  Theology. For the summary quotation see Tony Lane, A Concise History of Christian Thought.
  Revised. (Grand Rapids: Baker, 2006), 61.
\end{notelist}

\noteshead{Chapter 12}
\begin{notelist}
\entry St. Anthony’s experience in the desert may be found in Anne Fremantle, Treasury of Early
  Christianity (New York: Mentor, 1960), 400.
\entry The quotation describing Jerome is Roland Bainton’s from his Christendom, vol. 1, 135.
\entry The section “The Genius of the West” draws freely from Williston Walker’s sketch of Benedict
  in his Great Men of the Christian Church (Chicago: Chicago University Press, 1908), 103–14. The
  quotations in this section are from this work.
\end{notelist}

\noteshead{Chapter 13}
\begin{notelist}
\entry Augustine’s quotations about his early life are from his Confessions.
\entry The contrasts between Pelagius’s and Augustine’s views of sin and grace echo Williston
  Walker’s clear statement in Great Men of the Christian Church, 76–79.
\entry Augustine’s observation about the worldly city is from his City of God, especially book XIV,
  28 and book V, 19.
\entry On the Church’s reception of Augustine see Justo Gonzalez, The Story of Christianity: Volume
  One: The Early Church to the Dawn of the Reformation, revised and updated. (New York:
  HarperCollins, 2010), 250.
\end{notelist}

\noteshead{Chapter 14}
\begin{notelist}
\entry Valentinian’s decree and Leo’s quotation may both be found in Friedrich Gontard, The Chair of
  Peter (New York: Holt, Rinehart, and Winston, 1964), 138; 142–43.
\end{notelist}

\noteshead{Chapter 15}
\begin{notelist}
\entry The opening scene is from Kallistos Timothy Ware’s The Orthodox Church (Middlesex, Eng.:
  Penguin, 1964), 51.
\entry Harlie Kay Gallatin’s helpful description of icons is from Eerdman’s Handbook, 247–48.
\entry The words of the Russian envoys appear in Stephen Neill, Christian Missions (Middlesex:
  Penguin, 1964), 89.
\end{notelist}

\noteshead{Chapter 16}
\begin{notelist}
\entry Clovis’s prayer and baptism story appear in Roland Bainton, Christendom, vol. 1, 145–46.
\end{notelist}

\noteshead{Chapter 17}
\begin{notelist}
\entry Gregory the Great’s quotations in this chapter can be found in Philip Schaff, History of the
  Christian Church, vol. 4 (Grand Rapids: Eerdmans, 1950), 212–15, 228.
\entry Some biographical details in this chapter were taken from Robert Payne, The Fathers of the
  Western Church (New York: Viking, 1951).
\entry The story of the monk Copiosus is recorded in Odo J. Zimmerman, Saint Gregory: the Great
  Dialogues (New York: Fathers of the Church, 1959), 266–70.
\end{notelist}

\noteshead{Chapter 18}
\begin{notelist}
\entry The opening episode is reshaped from the account in Friedrich Gontard, The Chair of Peter,
  180–81.
\entry The politics of the Holy Roman Empire can be confusing to the layman. I have leaned heavily
  upon the treatment of a standard history text. In the sections “Ideas Die Slowly” and “Architect
  of an Empire” I have abridged Charlemagne’s rise to power from T. Walter Wallbank, Alastair M.
  Taylor, and Nels M. Bailkey, Civilization Past and Present (Glenview, IL: Scott,
\entry Foresman, 1975), 190–93. The explanation of feudalism is a condensation of pages 197–98, and
  the paragraphs dealing with the Germanic Kingdoms conflicts are an abbreviated version of pages
  243–45.
\entry David Bentley Hart, The Story of Christianity (London, Quercus: 2007),
\entry pp. 116–117.
\end{notelist}

\noteshead{Chapter 19}
\begin{notelist}
\entry Suger’s remark about the Gothic style appears in Anne Fremantle, Age of Faith (New York:
  Time-Life, 1965), 124.
\entry Pope Urban’s call to the crusade and the crusader quotation are both found in Civilization:
  Past and Present, 209–210.
\entry See Robert Wilken, The First One Thousand Years, (New Haven, CT:
\entry Yale, 2012), 288 and following for Christians under Muslim rule.
\end{notelist}

\noteshead{Chapter 20}
\begin{notelist}
\entry The scholar Olbert’s behavior is described in Civilization: Past and Present, 250.
\entry The description of “The Magnetism of an Able Teacher” comes from Age of Faith, 94–96. The
  quotations in this section are from this well-written work.
\end{notelist}

\noteshead{Chapter 21}
\begin{notelist}
\entry See the note on Albigensian Crusade by David Bentley Hart, The Story of Christianity (London:
  Quercus, 2007), 141.
\end{notelist}

\noteshead{Chapter 23}
\begin{notelist}
\entry John Hus’s prayer is quoted in Herbert B. Workman, Dawn of the Reformation, vol. 2 (London:
  Epworth, 1953), 325.
\entry For the inquisition note see Glenn Sunshine, The Reformation for Armchair Theologians
  (Louisville: Westminster John Knox Press, 2005), 96–99.
\end{notelist}

\noteshead{Chapter 24}
\begin{notelist}
\entry The best biography of Martin Luther is Roland Bainton, Here I Stand (Nashville: Abingdon,
  1950). It contains the quotations in this chapter, including the hymn describing Luther’s
  spiritual journey (66–67). In order to introduce these quotations, I have used the colorful
  summary of Luther’s life found in Time, March 24, 1967, 70–74, without attempting to identify
  quotations within quotation.
\entry For the theology and text boxes see Alister E. McGrath, Historical Theology: An Introduction
  to the History of Christian Thought, 2nd ed. (Malden, MA: Wiley-Blackwell, 2013).
\end{notelist}

\noteshead{Chapter 25}
\begin{notelist}
\entry The quotation from the martyred mother is found in Thielman J. Van Braght, Martyr’s Mirror
  (Scottdale, PA: Mennonite Publishing House, 1951), 984–87.
\entry I am indebted to John H. Yoder and Alan Kreider for their helpful summary of Anabaptist
  beliefs found in Eerdman’s Handbook, 399–403.
\end{notelist}

\noteshead{Chapter 28}
\begin{notelist}
\entry A brief summary of Loyola’s life and the Jesuit order appeared in Time, April 23, 1973,
  40–48. We have drawn some descriptions from this article.
\entry The quotations on the horrors of hell is from the Spiritual Exercises, 58.
\end{notelist}

\noteshead{Chapter 32}
\begin{notelist}
\entry The response to Ezra Stiles from Benjamin Franklin is quoted in Winthrop S. Hudson, American
  Protestantism (Chicago: University of Chicago Press, 1961), 13.
\entry The quotation from Baron von Holbach can be found in Civilization, 393–94.
\entry The Diderot quotation is found in Frank E. Manuel, The Age of Reason (Ithaca, NY: Cornell
  University Press, 1951), 30.
\entry For Renaissance Humanism see Alister E. McGrath. Historical Theology: An Introduction to the
  History of Christian Thought, 115–16.
\entry For the ancient wisdom of Revelation, see Robert Wilken, The Spirit of Early Christian
  Thought: Seeking the Face of God (New Haven, CT: Yale, 2003).
\end{notelist}

\noteshead{Chapter 33}
\begin{notelist}
\entry The quotations from the Pensées appear in many editions. Among them is the Modern Library
  edition, Blaise Pascal, Pensées (New York: Random House, 1941).
\entry The quote about August Herman Francke is as quoted by Tom Streeter in The Church and Western
  Culture: An Introduction to Church History (Bloomington, IN: AuthorHouse, 2006), 319.
\entry The sketch of Zinzendorf reflects Williston Walker’s treatment in Great Men of the Christian
  Church, 308–16.
\end{notelist}

\noteshead{Chapter 36}
\begin{notelist}
\entry The conditions and events of the French Revolution are condensed from Civilization, 451–61,
  and from Peter Gay, Age of Enlightenment (New York: Time-Life, 1966), 167–68.
\entry The account of the loss of the papal states and Vatican Council II is a summary of August
  Franzen’s treatments in A History of the Church, 384–94.
\end{notelist}

\noteshead{Chapter 37}
\begin{notelist}
\entry The quotation from G. M. Trevelyan can be found in Ernest Marshall Howse, Saints in Politics
  (London: George Allen, 1953), 178. I have found this book most helpful for this chapter.
\entry For the supplement to Bebbington’s description of Evangelicalism see Thomas Kidd, The Great
  Awakening; The Roots of Evangelical Christianity in Colonial America (New Haven, CT: Yale, 2007).
\end{notelist}

\noteshead{Chapter 38}
\begin{notelist}
\entry The attitude of the British East India Company is an echo of Hutchinson and Garrison, 20
  Centuries of Christianity, 279.
\entry The summary of Carey’s translation work is from Stephen Neill, The Christian Society (New
  York: Harper, 1952), 202.
\entry For a good treatment of Fuller see James Leo Garrett, Baptist Theology: a Four-Century Study.
  (Macon, GA: Mercer, 2007).
\entry I have tried to give voice to this immediacy in reading Scripture; see R. L. Hatchett. “The
  Hermeneutics of Conversion,” in The Ties That Bind: Life Together in the Baptist Vision, eds. Gary
  Furr and Curtis Freeman (Macon GA: Symth and Helwys, 1994).
\end{notelist}

\noteshead{Chapter 39}
\begin{notelist}
\entry The Timothy Flint portrait of the western preacher is quoted in Edwin Scott Gaustad,
  Historical Atlas of Religion in America (New York: Harper \& Row, 1962), 41.
\entry The Willard Sperry explanation of black religion was to an English audience: Religion in
  America (Boston: Beacon Press, 1963), 193.
\entry The quotations by President Woolsey and by H. L. Mencken come from George Marsden’s chapter
  in The Evangelicals, eds. David F. Wells and John D. Woodbridge (Nashville: Abingdon, 1975),
  122–23.
\end{notelist}

\noteshead{Chapter 40}
\begin{notelist}
\entry The summary of the “Spokesmen for Liberalism” draws upon Hordern, A Layman’s Guide, 44–49.
  The quotations are from these pages.
\end{notelist}

\noteshead{Chapter 41}
\begin{notelist}
\entry The Dickens quotations from Hard Times appears in Civilization: Past and Present, 487–88. The
  brief discussions of Marx and Engels are also drawn from this volume.
\entry The information on Bishop Ketteler and Cardinal Manning comes from Thomas Bokenkotter, A
  Concise History of the Catholic Church (Garden City, NY: Doubleday, 1977), 314–16.
\entry The Rauschenbusch explanation of the Social Gospel is from A Theology of the Social Gospel
  (New York: Macmillan, 1917), 5.
\end{notelist}

\noteshead{Chapter 42}
\begin{notelist}
\entry The discussion “Christians Under Hitler” draws upon Richard Pierard’s fine article, “An Age
  of Ideology,” found in the Eerdman’s Handbook, 576–78. The final quotation is from 587–88.
\end{notelist}

\noteshead{Chapter 43}
\begin{notelist}
\entry The quotations from Shailer Mathews and from J. Gresham Machen can be found in Robert L.
  Ferm, Issues in American Protestantism (Garden City, NY: Doubleday, 1969), 262–87.
\entry For the Scopes trial see Gary Wills, Under God: Religion and American Politics (New York:
  Simon and Schuster, 1900).
\end{notelist}

\noteshead{Chapter 44}
\begin{notelist}
\entry The George Whitefield quotation appears in William Warren Sweet, The Story of Religion in
  America (New York: Harper, 1950), 141–42.
\end{notelist}

\noteshead{Chapter 46}
\begin{notelist}
\entry See the article on “secularization” in The Blackwell Encyclopedia of Modern Christian Thought
  (Malden, MA: Blackwell, 1993).
\entry George Weigel, The Cube and the Cathedral: Europe, America, and Politics Without God (New
  York: Basic Books, 2005), 12.
\entry For smarter atheism see David Bentley Hart, Atheist Delusions: The Christian Revolution and
  its Fashionable Enemies (New Haven, CT: Yale, 2009) and Alister McGrath, The Twilight of Atheism:
  The Rise and Fall of Disbelief in the Modern World (WaterBrook Press, 2006).
\entry For stirrings among Muslims see Tom Doyle with Greg Webster, Dreams and Visions: Is Jesus
  Awakening the Muslim World? (Nashville, TN: Thomas Nelson, 2012).
\end{notelist}

\noteshead{Chapter 47}
\begin{notelist}
\entry Philip Jenkins’s “The Next Christianity” is found in The Atlantic, October, 2002, 53–68.
\entry Pentecostals stood in the vanguard of the Southern Counter-Reformation. Though Pentecostalism
  emerged as a movement only at the start of the twentieth century, chiefly in North America,
  Pentecostals a century later were at least 400 million strong and heavily concentrated in the
  global South. By 2040 or so there could be as many as a billion, at which point Pentecostal
  Christians alone would far outnumber the world’s Buddhists and would enjoy rough numerical parity
  with the world’s Hindus.
\entry For the roots and early history of Pentecostalism see Allan Anderson, An Introduction to
  Pentecostalism: Global Charismatic Christianity (New York: Cambridge, 2004), chapters 2 and 3.
\entry See Jenkins, The Next Christendom, 40 for the Kenyatta quote.
\entry For the ratio (1 in 4) of Charismatics to the total number of Christians see Pew Forum on
  Religion
\entry and Public Life (December 19, 2011), “Global Christianity: A Report on the Size and
  Distribution of
\entry the World’s Christian Population.”
\entry The chart on Pentecostal numbers is adapted from Mark Noll, The New Shape of World
  Christianity: How American Experience reflects Global Faith (Downers Grove, IL: InterVarsity
  Press, 2009), 22. Noll depends upon the numbers of David B. Barrett’s World Christian
  Encyclopedia, 2nd ed. (New York: Oxford University Press, 2001).
\entry See Noll, The New Shape of World Christianity,185 for East Africa.
\entry See A. Greg Carter, Rethinking Christ and Culture: A Post Christendom Perspective (Grand
  Rapids: Brazos Press, 2006) 14–15 for descriptions of Christendom.
\entry For Bunyan see Linda Woodhead, Introduction to Christianity (New York: Cambridge, 2004),
  204–5.
\entry See Mark Noll on the Bible and translation (relying upon Senneh), The Shape of World
  Christianity, 22–24, 36.
\entry For liberation theology see Paul R. Spickard and Kevin M. Cragg, A Global History of
  Christians: How Everyday Believers Experience Their World (Grand Rapids: Baker, 2004), 429.
\entry The “mass-produced affection” quotation is from Carol Flake, Redemptorama: Culture, Politics,
  and the New Evangelicalism (New York: Penguin Books, 1984), 17.
\entry In his Becoming Conversant with the Emerging Church (Grand Rapids: Zondervan, 2005), D. A.
  Carson writes, “At the heart of the Emergent Church movement—or as some of its leaders prefer to
  call it, the ‘conversation’—lies the conviction that changes in the culture signal that a new
  church is ‘emerging.’ Christian leaders must therefore adapt to this emerging church. Those who
  fail to do so are blind to the cultural accretions that hide the gospel behind forms of thought
  and modes of expression that no longer communicate with the new generation, the emerging
  generation” (p. 12). Many observers have associated the emerging church with postmodernity.
\entry Professor David E. Wells has written, “This postmodern outlook comes in all kinds of shapes
  and expressions, which is what probably explains the multiplicity of definitions which have been
  advanced. Its own ethos almost guarantees that there will be no such thing as a postmodern outlook
  but rather there will be many different postmodern perspectives. Yet what they have in common is
  that they all believe that meaning has died. This has to be qualified immediately by the assertion
  that what has most obviously died is the kind of rational meaning which the Enlightenment
  provided—but postmoderns find no grounds for any other kind.” From Above All Earthly Pow’rs (Grand
  Rapids: Eerdmans, 2005), 67.
\end{notelist}

\noteshead{Chapter 48}
\begin{notelist}
\entry The description of the underground church scene is from “China’s Christian Underground,” by
  Fang Bay, U.S. News \& World Report, April 30, 2001. And the striking shift in the global map is
  found in Christopher J. H. Wright’s “An Upside Down World” in Christianity Today, January 2007,
  42.
\entry David Aikman’s story and report from China can be found in his Jesus in Beijing. The Su
  Wenxing story is on pages 258–60.
\entry In the early years of the twenty-first century, the Christian population of South Korea also
  continued at a rapid pace of nearly 6 percent annually, bolstered by vigorous evangelistic efforts
  and also by the favor of the government, which viewed Christianity as an ideological means of
  resisting the encroachment of communism.
\entry For the stories I depend upon Mark Noll and Carolyn Nystrom, Clouds of Witnesses: Christian
  Voices from Africa and Asia (Downer Grove, IL: InterVarsity Press, 2011); Jonathan Hill, Zondervan
  Handbook to the History of Christianity (Oxford: Zondervan, 2006); Allen Anderson, To the Ends of
  the Earth: Pentecostalism and the Transformation of World Christianity (New York: Oxford
  University Press, 2013); Paul R. Spickard and Kevin M. Cragg, A Global History of Christians: How
  Everyday Believers Experience Their World (Grand Rapids: Baker, 2004); and Douglas Jacobsen, The
  World’s Christians: Who they Are, Where They Are, and How They Got There (Oxford, UK:
  Wiley-Blackwell, 2011).
\entry The Carter story is related by Daniel Vestal; cf. also
  \url{http://www.washington-post.com/wp-srv/newsweek/religion.htm}
\entry For the conservative numbers I depended upon the Pew Forum on Religion and Public Life
  (December 19, 2011), “Global Christianity: A Report on the Size and Distribution of the World’s
  Christian Population,” 67.
\entry For the East African revival I follow Mark Noll, The New Shape of World Christianity,
  169–183.
\entry On line, The Dictionary of African Christian Biography provides well-informed articles;
  numerous other sources are helpful including Wesley L. Handy’s \url{missionforum.wordpress.com}.
\end{notelist}

\noteshead{For Further Study}
\begin{notelist}
\entry In addition to the books mentioned at the end of each chapter and in these notes, other
  reference books offer leads to material to satisfy the curiosity of the most eager student. The
  Oxford Dictionary of the Christian Church, 3rd ed., revised, edited by F. L. Cross and Elizabeth
  A. Livingstone (New York: Oxford University Press, 2005) and The New International Dictionary of
  the Christian Church, edited by J. D. Douglas (Grand Rapids: Zondervan, 1978) have articles and
  additional suggestions for reading on the issues within the chapters of this volume.
\entry To pursue some interest stimulated by the chapters in this volume dealing with American
  Christianity, the student can also find articles and book suggestions in Dictionary of
  Christianity in America, edited by Daniel C. Reid (Downers Grove, IL: InterVarsity, 1990). A whole
  range of maps for the study of Christian history appear in The Macmillan Atlas History of
  Christianity by Franklin H. Littell (New York: Macmillan, 1976).
\end{notelist}
